% !TeX program = xelatex
\documentclass[11pt,a4paper]{article}
\usepackage[margin=24mm]{geometry}
\usepackage{fontspec,xeCJK}
\setmainfont{lmroman10-regular.otf}[BoldFont=lmroman10-bold.otf,ItalicFont=lmroman10-italic.otf]
\setCJKmainfont{FandolSong-Regular.otf}[BoldFont=FandolSong-Bold.otf,ItalicFont=FandolKai-Regular.otf]
\usepackage{amsmath,booktabs,graphicx,tikz}
\usetikzlibrary{positioning,arrows.meta}
\usepackage[unicode,hidelinks]{hyperref}
\renewcommand{\abstractname}{摘要占位}
\renewcommand{\refname}{参考文献}
\renewcommand{\figurename}{图}
\renewcommand{\tablename}{表}
\setlength{\parindent}{2em}\setlength{\parskip}{0.3em}\sloppy
\title{不同训练样本规模下空间注意力对焊缝X射线图像\\多类别无监督聚类的影响}
\author{作者与单位待填写}
\date{研究骨架 v0.5（未冻结）\quad 2026年11月27日}
\begin{document}\maketitle
\begin{abstract}
\textbf{目的：}研究目标域类别标签不参与表征训练和聚类拟合时，空间注意力对无缺陷及不同缺陷类型图像的聚类影响，以及这种影响随训练样本规模的变化。
\textbf{数据审计：}本次接收32,932张TIFF和13,859份XML，完成全量结构检查及独立复算；最终科研样本池、类别和来源组资格尚未确认。
\textbf{方法占位：}候选框架为共享训练设置的ResNet18与SimCLR，对比有、无空间分支；低、中、高规模、K值、插入位置及训练预算待数据资格和训练侧Pilot后冻结。
\textbf{结果占位：}尚无模型训练或聚类结果。
\textbf{结论占位：}待依据冻结协议的真实实验填写。
\end{abstract}

\section{引言}
本研究考察空间加权是否改变无缺陷与不同缺陷类型的表征及聚类，并将无标签训练规模作为实验因素。焊缝监督二分类中的注意力研究\cite{wang2025attention}与工业异常类型的局部加权聚类\cite{sohn2023anomaly}构成相关背景。论文贡献待实验和文献排重后确定，不将既有模块的组合当作新注意力算法。

\section{数据与研究资格}
\subsection{已执行的结构审计}
2026年11月27日接收目录中包含32,932张TIFF，均能完整解码，为16位单通道单帧。缺陷目录含13,859张，用户说明合格目录中的19,073张代表合格图像；后者没有XML，合格身份属于研究者说明，未通过本次专业判读认证。13,859份XML均可解析，其宽高与对应图像一致。标注对象共31,898个，包括焊缝辅助区域。

检测到255组完全重复内容，合计256个多余文件，其中17组跨现成分区；25组相同像素具有不同XML缺陷标签集合。目录类别未在XML缺陷集合出现的图像为1,199张；具有多个XML缺陷类型的图像为1,899张；218张XML只有焊缝辅助标签。上述集合可能重叠。它们是待复核事项，不能直接解释为已知真值下的错标率。

图像共有1,666种宽高组合；没有原始底片及导出记录，不能仅由尺寸认证逐张裁剪历史。47张不同图像完成AI视觉分诊，未替代无损检测专家逐张核验。结构候选数、像素唯一数均不是独立焊口数。审计方法与机器汇总随仓库文档提供，逐文件证据及工业图像不公开。

\subsection{拟议评价范围}
观察单位为一张提供的焊缝区域图像；拟研究无缺陷与经专业确认的不同缺陷类型。多缺陷图像保留多标签事实，单标签主评价资格待定。其它缺陷为含义待澄清的集合标签，不自动作为同质语义类别。无缺陷及具名缺陷候选均需标注方确认。

同一焊口、源图或重复连通组件不得跨分区。当前来源证据不足，数据资格保持HOLD，未生成最终划分。本次为预分区审计，已访问现成train/val/test目录标签；这些目录不构成尚未接触的最终测试集。后续冻结研究数据版本、来源组划分、独立Pilot协议和最终分析规范。

\section{候选方法}
Baseline候选为ResNet18与SimCLR\cite{he2016resnet,chen2020simclr}，Attention候选在相同训练流程中加入空间分支。若使用CBAM空间分支\cite{woo2018cbam}，应说明其区别于完整CBAM；插入位置待训练侧Pilot决定。若使用ImageNet监督预训练，披露其外部监督来源。目标域类别标签不用于表征训练、主训练子集选择、模型选择或聚类拟合。

K-means为候选算法。K值规则未确定，不能继承二分类K=2。若使用已知类别数，披露该先验；若研究未知K，在独立无标签拟合池上预先规定选K规则。最终测试不参与PCA、聚类拟合或参数选择。

\begin{figure}[ht]\centering
\begin{tikzpicture}[node distance=6mm and 7mm,b/.style={draw,rounded corners,align=center,text width=27mm,minimum height=15mm,font=\small},a/.style={-{Stealth[length=2mm]},thick}]
\node[b](d){全量结构审计\\专业语义复核};
\node[b,right=of d](s){来源组与重复\\资格确认};
\node[b,right=of s](t){训练侧Pilot\\最终协议冻结};
\node[b,right=of t](m){低/中/高规模\\配对模型};
\node[b,below=of m](f){冻结表示\\聚类拟合集};
\node[b,left=of f](e){独立测试评价\\性能/几何/稳定性};
\draw[a](d)--(s);\draw[a](s)--(t);\draw[a](t)--(m);\draw[a](m)--(f);\draw[a](f)--(e);
\end{tikzpicture}
\caption{拟议流程；当前完成结构审计，专业语义与来源资格未通过。}
\end{figure}

\section{实验与分析计划}
比较低、中、高三档无标签训练规模下的配对模型，实际数量、预算与重复次数待资格及真实吞吐测量确定。对各规模记录唯一图像数、独立组数、重复曝光、墙钟时间与内存。固定更新步数与固定epoch的解释不同，预算规则在冻结前明确。

ARI为外部指标候选，AMI、NMI为补充候选\cite{vinh2010information}。整体多类别评价与在同一簇分配上对缺陷子集的诊断分别报告。类内/类间几何、kNN purity与稳定性为解释性分析候选；其定义、覆盖与无效规则在最终评价前冻结。注意力区域关联需有可信坐标与语义标签；热图和几何变化不能单独建立因果机制。

令$\Delta_s(r)=ARI_{\mathrm{attention},s}(r)-ARI_{\mathrm{baseline},s}(r)$。逐规模报告原始结果和配对差，低高规模的$\Delta$差作为规模作用比较候选；主要对比、统计单位和汇总尚未冻结，不根据最终测试结果选择。

\section{模型结果占位}
\begin{table}[ht]\centering
\caption{模型结果尚未产生；横线表示未实验。}
\begin{tabular}{lccc}\toprule
训练规模&Baseline ARI&Attention ARI&配对差\\\midrule
低&---&---&---\\中&---&---&---\\高&---&---&---\\\bottomrule
\end{tabular}\end{table}
待填聚类、几何、稳定性、覆盖及全部不利或无明显差异的结果。数据结构审计数字不填入模型性能表。

\section{讨论与结论占位}
待依据真实实验回答空间注意力的作用及其规模依赖，讨论缺陷信息与来源线索等竞争性解释。当前只有结构审计与有限AI视觉分诊，不判断注意力有效或无效，不推断完整因果机制。
\bibliographystyle{plain}\bibliography{references}
\end{document}
